\begin{prob}
    In each of the problems below state the best case and worst case time
    complexities of the given piece of code using asymptotic notation. Note that
    some algorithms may have the same best case and worst case time complexities. If
    the best and worst case complexities \emph{are} different, identify which
    inputs result in the best case and worst case. You do not otherwise need to show your work for
    this problem.

    \emph{Example Algorithm:} \python{linear_search} as given in lecture. 

    \emph{Example Solution:}
    Best case: $\Theta(1)$, when the target is the first element of the array.
    Worst case: $\Theta(n)$, when the target is not in the array.

    \begin{subprobset}
        \begin{subprob}
            \begin{minted}[autogobble]{python}
                def first_repeated(data):
                    """Finds the first repeated element in `data`.
                    Returns `None` if no elements are repeated.
                    `data` is an array of size n."""
                    n = len(data)
                    for i in range(n):
                        for j in range(i + 1, n):
                            if data[i] == data[j]:
                                return data[i]
                    return None
            \end{minted}

            \begin{soln}
                Best case: $\Theta(1)$, when the first and second elements of the array
                are the same.

                Worst case: $\Theta(n^2)$, when there are no repeated elements.

                You have to wonder... isn't there a faster approach?
            \end{soln}
        \end{subprob}

        \begin{subprob}
            \begin{minted}[autogobble]{python}
                def f_2(data):
                    """`data` is an array of n numbers"""
                    m = data[0]
                    c = 1
                    for ix in range(1, len(data)):
                        if c == 0:
                            m = data[ix]
                            c = 1
                        elif m == data[ix]:
                            c += 1
                            if c > len(data) // 2:
                                return m
                        else:
                            c -= 1

                    return m
            \end{minted}
            

            \begin{soln}
                Best case and worst case are both $\Theta(n)$. In the ``best case'',
                the loop exits early when \python{c > len(data) // 2}. But $c$ grows by
                at most 1 on each iteration, so we'd need to go through at least
                $n / 2$ iterations to exit, which is still $\Theta(n)$.

                By the way, this algorithm is called the Boyer-Moore majority algorithm,
                and can be used to efficiently find a majority element of a list (as long
                as there is one). Though note that you can reason about the
                algorithm's time complexity without actually knowing what it
                computes.
            \end{soln}
        \end{subprob}

        \begin{subprob}
            \begin{minted}[autogobble]{python}
                def kth_smallest(numbers, k):
                    """Finds the k-th smallest element in the array.
                    `numbers` is an array of n numbers."""
                    n = len(numbers)
                    for i in range(n):
                        count = 0 # Count how many numbers are less than numbers[i]
                        for j in range(n):
                            if numbers[j] < numbers[i]:
                                count += 1
                        if count == k - 1:
                            return numbers[i]
            \end{minted}
            
            \begin{soln}
                Best case: $\Theta(n)$, when the first element of the array is the $k$th smallest.

                Worst case: $\Theta(n^2)$, we need to count the number of elements smaller than each element. If the last element is $k$th smallest, then we need to run entire nested loops.
            \end{soln}
        \end{subprob}

        \begin{subprob}
            \begin{minted}[autogobble]{python}
                def index_of_kth_smallest(numbers, k):
                    """`numbers` is an array of size n"""
                    # the kth_smallest() from above
                    element = kth_smallest(numbers)
                    # the linear_search() from lecture
                    return linear_search(numbers, element)
            \end{minted}
            

            \begin{soln}
                Best case: $\Theta(n)$. This occurs when the $k$th smallest element is the first element of the array,
                as then \python{kth_smallest} takes $\Theta(n)$ and \python{linear_search} takes $\Theta(1)$,
                for a total of $\Theta(n)$.

                Worst case: $\Theta(n^2)$. This occurs when the $k$th smallest element is the last element of the
                array. In this situation, \python{kth_smallest} takes $\Theta(n^2)$ time and \python{linear_search}
                takes $\Theta(n)$ time, for a total of $\Theta(n^2)$ time.
            \end{soln}
        \end{subprob}

    \end{subprobset}


\end{prob}
